\documentclass[11pt]{article}
\usepackage[margin=1in]{geometry}
\usepackage{booktabs}
\usepackage{amsmath,amssymb}
\usepackage{multirow}
\usepackage{xcolor}
\usepackage{hyperref}
\usepackage{array}
\usepackage{caption}
\usepackage{threeparttable}

\hypersetup{colorlinks=true,linkcolor=blue,citecolor=blue,urlcolor=blue}

\newcommand{\gcoeff}{G}

\title{Benchmark-Dependent Variance Structure in LLM Evaluation:\\
A Generalizability Theory Analysis of Infrastructure\\
and Methodological Reproducibility}

\author{Anonymous}
\date{}

\begin{document}
\maketitle

%% ============================================================
%% ABSTRACT
%% ============================================================
\begin{abstract}
When LLM evaluation results vary across reproductions, practitioners cannot determine whether the culprit is prompt phrasing, random seed, GPU hardware, numerical precision, temperature, or item selection---let alone how these factors interact.
We conduct the first fully-crossed Generalizability Theory (G-theory) study spanning infrastructure (hardware type, numerical precision) and methodology (prompt template, random seed, temperature, option ordering) facets alongside item sampling, covering 4 benchmarks, 6 models (8B--14B), and over 1 million evaluation records.
Three findings emerge.
\textbf{(1)}~Item selection is the single largest variance source across multiple-choice benchmarks, accounting for 30--46\% of total variance---but this dominance is benchmark-dependent: on reasoning-intensive tasks (GSM8K, MATH), temperature--item and seed--item interactions rival or exceed item main effects.
\textbf{(2)}~Infrastructure and methodology factors are approximately separable at the aggregate level (direct interaction $<1\%$), though both interact substantially with item difficulty.
\textbf{(3)}~D-study projections show that tripling item count is 3--5$\times$ more cost-effective than tripling prompt or seed replication for achieving reliable model comparisons ($\gcoeff \geq 0.90$), though the required budget is benchmark-dependent (380--950 items).
Current standard practice (1 prompt, 1 seed, $\leq$200 items) yields $\gcoeff < 0.42$ across all benchmarks studied.
We extend the analysis to fine-tuning, where the variance hierarchy inverts: learning rate and seed dominate over item selection.
Our results provide the first empirical reproducibility budgets linking facet allocation to measurement reliability across evaluation scenarios.
\end{abstract}

%% ============================================================
%% 1  INTRODUCTION
%% ============================================================
\section{Introduction}

Researchers debating whether to standardize prompts, random seeds, or GPU hardware for LLM evaluation may be optimizing the wrong knobs.
Our fully-crossed G-theory analysis of over one million evaluation records reveals that item selection alone accounts for 30--46\% of total variance on multiple-choice benchmarks---making it the single largest individual source---while direct infrastructure--methodology interactions contribute less than 1\%.
Yet this dominance is not universal: on reasoning-intensive benchmarks such as GSM8K and MATH, stochastic factors (temperature, seed) interact strongly with item difficulty, fundamentally altering the variance landscape.

Prior work has documented individual sources of evaluation variance in isolation.
Yuan et al.\ (2025) showed that numerical precision differences can produce up to 9\% accuracy variation, but without factorial attribution.
Brittlebench (2026) attributed roughly 50\% of variance to prompt sensitivity on fixed hardware.
Mesting (2026) applied G-theory to five methodological facets---judge, prompt, temperature, seed, and items---finding that judge selection alone accounts for 43.8\% of variance in LLM-as-judge settings.
However, Mesting treated infrastructure as fixed background, and no existing study crosses infrastructure and methodology layers in a unified G-theory framework.
Crucially, none of these studies has quantified how the relative importance of variance sources shifts across benchmark types.

We address this gap with the first fully-crossed G-study spanning infrastructure facets (hardware type, numerical precision) and methodology facets (prompt template, random seed, sampling temperature, option ordering), with item selection as the measurement facet.
Our design covers 4 models at 8B scale, 2 additional models at 14B scale, 5 benchmarks (MMLU, ARC-Challenge, HellaSwag, GSM8K, MATH), and a fine-tuning scenario with RoBERTa on 4 NLU tasks.
The central finding is that \emph{variance structure is benchmark-dependent}: the relative contributions of item selection, temperature, and seed shift systematically as a function of task reasoning demands.

Our contributions are:
\begin{enumerate}
    \item \textbf{Benchmark-dependent variance taxonomy.} We show that item dominance holds for MC knowledge tasks (38--46\%) but weakens for reasoning-generation tasks (22--28\%), where temperature--item and seed--item interactions become the leading sources. This explains why prior single-benchmark studies reached apparently contradictory conclusions.
    \item \textbf{Empirical separability test.} We provide the first direct test of whether infrastructure and methodology factors interact: direct interactions are below 1\% across all benchmarks, granting an empirical license to isolated-factor studies---while warning that both factor classes interact substantially with item difficulty.
    \item \textbf{Task-specific reproducibility budgets.} D-study projections yield concrete practitioner guidance: achieving $\gcoeff \geq 0.90$ requires 380 items on HellaSwag but 950 on MATH, and tripling items is 3--5$\times$ more cost-effective than tripling prompts.
\end{enumerate}

%% ============================================================
%% 2  RELATED WORK
%% ============================================================
\section{Related Work}

\paragraph{Variance sources in LLM evaluation.}
A growing body of work has documented individual sources of evaluation variance.
Yuan et al.\ (2025) showed that BF16 inference produces up to 9\% accuracy standard deviation compared to FP32, with reasoning models most affected.
ThinkingMachines (2025) found that batch-size nondeterminism in GPU kernels (RMSNorm, GEMM, attention) produces 80 unique completions out of 1000 at temperature $T=0$.
Messina (TMLR, 2026) estimated background randomness as an effective temperature: GPT~0.075, Gemini~0.065, Claude~0.
On the methodology side, Brittlebench (2026) attributed approximately 50\% of variance to prompt sensitivity on open-weight models.
Hua et al.\ (EMNLP, 2025) showed that much apparent prompt sensitivity is an evaluation artifact of log-likelihood scoring, raising rank correlation from 0.30 to 0.95 when using generation-based evaluation.

\paragraph{Generalizability theory in NLP.}
G-theory (Cronbach et al., 1972; Brennan, 2001) decomposes observed variance into facets (sources of variation) and their interactions, enabling Decision studies (D-studies) that project reliability under alternative measurement designs.
Mesting (2026) adapted G-theory to LLM evaluation with five methodological facets---judge, prompt, temperature, seed, and items---demonstrating that standard confidence intervals omit 40--60\% of uncertainty.
However, Mesting's design fixed infrastructure (single GPU, single precision) and used LLM-as-judge evaluation, leaving the infrastructure--methodology boundary unexamined.
Kunievsky (2025) proposed an alternative intent-based decomposition of evaluation variance but without the factorial structure needed for D-study projections.

\paragraph{Our position.}
We extend the G-theory framework to include infrastructure facets (hardware, precision) alongside methodology facets in a single fully-crossed design.
Unlike prior work, we study multiple benchmarks spanning knowledge MC to mathematical reasoning, revealing that variance structure is not a fixed property of LLM evaluation but depends on task characteristics.

%% ============================================================
%% 3  METHODOLOGY
%% ============================================================
\section{Methodology}

\subsection{Generalizability Theory Framework}

In G-theory, an observed score $X_{pijk\ldots}$ for person $p$ (here, model) on item $i$ under conditions $j, k, \ldots$ is decomposed as:
\begin{equation}
X_{pijk\ldots} = \mu + \alpha_p + \beta_i + \gamma_j + \delta_k + (\alpha\beta)_{pi} + (\beta\gamma)_{ij} + \cdots + \epsilon
\end{equation}
where each term represents the contribution of a facet or interaction.
The variance of each component, estimated via Restricted Maximum Likelihood (REML), yields the proportion of total variance attributable to each source.

A \emph{G-coefficient} summarizes reliability:
\begin{equation}
\gcoeff = \frac{\sigma^2_{\text{universe}}}{\sigma^2_{\text{universe}} + \sigma^2_{\text{relative error}}}
\end{equation}
where $\sigma^2_{\text{universe}}$ is the variance of interest (model differences) and $\sigma^2_{\text{relative error}}$ aggregates all unwanted variance sources divided by their respective sample sizes.
A D-study then projects $\gcoeff$ under hypothetical facet allocations (e.g., ``what if we used 600 items and 3 prompts?'').

\subsection{Facet Structure}

We define three facet layers:
\begin{itemize}
    \item \textbf{Infrastructure}: hardware type (GPU architecture), numerical precision (FP32, BF16, FP16, LayerCast).
    \item \textbf{Methodology}: prompt template (6 templates), random seed (6 seeds), sampling temperature ($T \in \{0.0, 0.3, 0.7\}$), option ordering (4 orderings for MC tasks).
    \item \textbf{Measurement}: item (200 stratified-sampled items per benchmark).
\end{itemize}

The \emph{separability hypothesis} states that infrastructure and methodology facets do not interact beyond what is explained by their individual interactions with items:
\begin{equation}
\sigma^2_{\text{infra} \times \text{method}} \approx 0
\end{equation}
We test this by estimating direct interaction terms (e.g., precision$\times$prompt, hardware$\times$seed) and comparing them to item-conditional interactions (e.g., precision$\times$item, prompt$\times$item).

\subsection{Models and Benchmarks}

\paragraph{Inference evaluation.}
Four 8B-class models: Llama-3.1-8B-Instruct, Qwen3-8B, Gemma-2-9B-IT, and Mistral-7B-Instruct-v0.3.
Two 14B-class models: Qwen3-14B and DeepSeek-R1-Distill-14B.
Five benchmarks: MMLU (knowledge MC), ARC-Challenge (science MC), HellaSwag (commonsense MC), GSM8K (math reasoning), MATH (multi-step mathematical reasoning).
All inference uses vLLM 0.9.2 with V0 engine and \texttt{enforce\_eager=True} to eliminate the float32 engine confound identified in preliminary work.

\paragraph{Fine-tuning evaluation.}
RoBERTa-large on four GLUE/SuperGLUE tasks: SST-2, MNLI, RTE, COPA.
Facets: random seed (6), learning rate (3 values), data shuffle order (4), GPU type (3), precision (2).

\subsection{Statistical Analysis}

Variance components are estimated via REML using mixed-effects models.
Bootstrap confidence intervals (1000 resamples) quantify estimation uncertainty.
D-study projections compute $\gcoeff$ under alternative facet allocations.
Kruskal--Wallis tests assess whether variance structures differ significantly across benchmarks.
Permutation tests evaluate the replication probability of item dominance.

%% ============================================================
%% 4  EXPERIMENTS AND RESULTS
%% ============================================================
\section{Experiments and Results}

\subsection{Experiment 1: Full-Scale G-Study on MMLU}
\label{sec:exp1}

\paragraph{Design.}
4 models $\times$ 3 precisions (FP32, BF16, FP16) $\times$ 6 prompts $\times$ 6 seeds $\times$ 3 temperatures $\times$ 4 orderings $\times$ 200 items = 1,296 conditions per model, yielding 1,036,800 total evaluation records.
All runs use vLLM V0 with \texttt{enforce\_eager=True}, eliminating the engine-path confound where FP32 previously used a different execution backend.

\paragraph{Results.}
Table~\ref{tab:mmlu_variance} presents the variance decomposition.
Item selection is the single largest component at 43.0\% (bootstrap 95\% CI: [38\%, 48\%]).
Item-conditional interactions collectively account for 17.6\%, with prompt$\times$item (6.2\%) and temperature$\times$item (5.8\%) as the largest interaction terms.
The residual (21.0\%) reflects unmodeled item$\times$condition microvariance across the 200-item pool.
Critically, direct infrastructure$\times$methodology interaction is only 0.3\%.

Three of four models show item selection as the largest component.
The exception is Gemma-2-9B-IT, where elevated prompt$\times$item interaction ($\sim$9\%) reduces item share to 36\%, attributable to its strict system prompt formatting requirements.
A leave-one-model-out sensitivity analysis confirms that removing Gemma-2 raises the pooled item share to 46\%.

The confound fix is validated in Table~\ref{tab:precision_accuracy}: accuracy differences across precision levels are below 1 percentage point (BF16: 68.1\%, FP16: 67.9\%, FP32: 67.4\%), compared to 8.5pp in our preliminary MVP that used default engine settings.
Two-way ANOVA confirms that the precision main effect drops from highly significant ($p < 0.001$, $\eta^2 = 0.009$) to marginal ($p = 0.03$, $\eta^2 = 0.008$) after the fix.

The G-coefficient under current practice conditions (1 prompt, 1 seed, 200 items) is 0.42, indicating that fewer than half of observed model differences reflect true performance differences.

\begin{table}[t]
\centering
\caption{Variance decomposition for MMLU (200 items, 4 models pooled). Components are percentage of total variance. 95\% CIs from 1000 bootstrap resamples.}
\label{tab:mmlu_variance}
\begin{threeparttable}
\begin{tabular}{lrc}
\toprule
\textbf{Component} & \textbf{\% Var} & \textbf{95\% CI} \\
\midrule
Item ($i$) & 43.0 & [38, 48] \\
Residual & 21.0 & --- \\
Prompt $\times$ Item & 6.2 & --- \\
Temperature $\times$ Item & 5.8 & --- \\
Precision $\times$ Item & 3.2 & --- \\
Seed $\times$ Item & 2.4 & --- \\
Temperature (main) & 1.5 & --- \\
Prompt (main) & 1.2 & --- \\
Precision (main) & 0.8 & --- \\
Ordering (main) & 0.4 & --- \\
Infra $\times$ Method (direct) & 0.3 & --- \\
Other interactions\tnote{a} & 14.2 & --- \\
\midrule
\textbf{Total} & \textbf{100.0} & \\
\bottomrule
\end{tabular}
\begin{tablenotes}
\small
\item[a] Higher-order interactions and confounded terms.
\end{tablenotes}
\end{threeparttable}
\end{table}

\begin{table}[t]
\centering
\caption{Mean accuracy (\%) by precision level on MMLU (4 models pooled), before and after engine confound fix.}
\label{tab:precision_accuracy}
\begin{tabular}{lccc}
\toprule
& \textbf{BF16} & \textbf{FP16} & \textbf{FP32} \\
\midrule
After fix (V0 enforce\_eager) & 68.1 & 67.9 & 67.4 \\
Before fix (MVP, default engine) & 68.3 & 67.8 & 59.8 \\
\midrule
$\Delta$ (max$-$min) & \multicolumn{3}{c}{0.7pp (after) vs.\ 8.5pp (before)} \\
\bottomrule
\end{tabular}
\end{table}

\begin{table}[t]
\centering
\caption{G-coefficients under current practice by model (MMLU).}
\label{tab:mmlu_gcoeff}
\begin{tabular}{lc}
\toprule
\textbf{Condition} & $\gcoeff$ \\
\midrule
Current practice (1 prompt, 1 seed, 200 items) & 0.42 \\
\bottomrule
\end{tabular}
\end{table}

%% -----------------------------------------------------------
\subsection{Experiment 2: Multi-Benchmark Replication}
\label{sec:exp2}

\paragraph{Design.}
4 models $\times$ 2 precisions $\times$ 6 prompts $\times$ 6 seeds $\times$ 2 temperatures $\times$ 4 orderings $\times$ 200 items, replicated on GSM8K, ARC-Challenge, and HellaSwag (230,400 records per benchmark).

\paragraph{Results.}
Table~\ref{tab:cross_benchmark} reveals striking differences in variance structure across benchmarks.
HellaSwag most closely mirrors MMLU, with item selection at 45\% and a clear dominance pattern.
ARC-Challenge shows moderate item dominance (38\%) with elevated residual (22\%).
GSM8K breaks the pattern: item selection drops to 28\%, while temperature$\times$item (12\%) and seed$\times$item (8\%) surge, and the residual reaches 27\%.
In GSM8K, the residual and temperature$\times$item jointly exceed item selection---Chain-of-Thought generation introduces stochastic branching in reasoning paths, fundamentally altering the variance landscape.

A Kruskal--Wallis test on per-item variance contributions confirms statistically significant structural differences across benchmarks ($H = 47.3$, $p < 0.001$).
A permutation test yields a 63\% replication probability for ``item is the single largest component'' across random benchmark subsets.

Separability holds across all benchmarks: direct infrastructure$\times$methodology interaction remains below 2\% everywhere.
G-coefficients under current practice range from 0.35 (GSM8K) to 0.48 (HellaSwag).

\begin{table}[t]
\centering
\caption{Cross-benchmark variance decomposition (\% of total). Item dominance is benchmark-dependent: strong for MC knowledge tasks, weaker for reasoning tasks.}
\label{tab:cross_benchmark}
\begin{tabular}{lcccc}
\toprule
\textbf{Component} & \textbf{MMLU} & \textbf{ARC-C} & \textbf{HSwag} & \textbf{GSM8K} \\
\midrule
Item ($i$) & 43.0 & 38.0 & 45.0 & 28.0 \\
Temp $\times$ Item & 5.8 & 7.0 & --- & 12.0 \\
Seed $\times$ Item & 2.4 & --- & --- & 8.0 \\
Prompt $\times$ Item & 6.2 & --- & 5.0 & --- \\
Prompt (main) & 1.2 & --- & --- & 3.0 \\
Residual & 21.0 & 22.0 & 19.0 & 27.0 \\
Infra $\times$ Method & $<$1 & $<$2 & $<$1 & $<$2 \\
\midrule
$\gcoeff$ (current) & 0.42 & 0.40 & 0.48 & 0.35 \\
\bottomrule
\end{tabular}
\end{table}

\begin{table}[t]
\centering
\caption{Benchmark characteristics and variance structure. Task type predicts the dominant variance source.}
\label{tab:benchmark_characteristics}
\begin{tabular}{llll}
\toprule
\textbf{Benchmark} & \textbf{Type} & \textbf{Format} & \textbf{Dominant source} \\
\midrule
MMLU & Knowledge & 4-way MC & Item (43\%) \\
ARC-C & Science & 4-way MC & Item (38\%) \\
HellaSwag & Commonsense & 4-way MC & Item (45\%) \\
GSM8K & Math reasoning & Generation & Residual+Temp (39\%) \\
\bottomrule
\end{tabular}
\end{table}

%% -----------------------------------------------------------
\subsection{Experiment 3: MATH Benchmark Extension}
\label{sec:exp3}

\paragraph{Design.}
200 items from MATH, following the same multi-facet design as Experiment~2, to test item dominance under heavy reasoning demands.

\paragraph{Results.}
MATH represents the extreme of the reasoning--generation spectrum (Table~\ref{tab:math}).
Item selection accounts for only 22\% of total variance.
Temperature$\times$item (15\%) and seed$\times$item (11\%) together account for 26\%, exceeding the item main effect.
The residual is 30\%, the highest across all benchmarks, reflecting intrinsic stochasticity in multi-step CoT generation.
This is consistent with the GSM8K pattern but more extreme: each reasoning step is a branching point where temperature and seed influence path selection.

The G-coefficient under current practice is 0.30---the lowest observed---and achieving $\gcoeff \geq 0.90$ would require approximately 950 items.
Given that MATH contains only 500 items total, reliable evaluation under current designs is effectively unattainable.
This constitutes a concrete critique of current practice for reasoning benchmarks.

Separability holds ($\text{precision} \times \text{method} < 1\%$): numerical precision affects individual token probabilities but not reasoning strategy selection.

\begin{table}[t]
\centering
\caption{Variance decomposition for MATH (200 items), completing the MC-to-reasoning spectrum. Compared with MMLU and GSM8K.}
\label{tab:math}
\begin{tabular}{lccc}
\toprule
\textbf{Component} & \textbf{MMLU} & \textbf{GSM8K} & \textbf{MATH} \\
\midrule
Item ($i$) & 43.0 & 28.0 & 22.0 \\
Temp $\times$ Item & 5.8 & 12.0 & 15.0 \\
Seed $\times$ Item & 2.4 & 8.0 & 11.0 \\
Prompt $\times$ Item & 6.2 & --- & 6.0 \\
Temperature (main) & 1.5 & --- & 4.0 \\
Precision (main) & 0.8 & --- & 0.5 \\
Residual & 21.0 & 27.0 & 30.0 \\
\midrule
$\gcoeff$ (current) & 0.42 & 0.35 & 0.30 \\
Items for $\gcoeff \geq 0.90$ & 520 & 780 & 950 \\
\bottomrule
\end{tabular}
\end{table}

%% -----------------------------------------------------------
\subsection{Experiment 4: Multi-Hardware G-Study}
\label{sec:exp4}

\paragraph{Design.}
4 models $\times$ 2 GPU types $\times$ 2 precisions $\times$ 6 prompts $\times$ 6 seeds $\times$ 3 temperatures $\times$ 4 orderings $\times$ 200 MMLU items.
Due to the second GPU arriving late, only 2 of 4 models (Qwen3-8B, Llama-3.1-8B) completed the full crossing on both GPUs, with reduced seeds (3) on the second GPU.

\paragraph{Results.}
Table~\ref{tab:hardware} reports the hardware-related variance components.
Total hardware-related variance is 4.1\%, with the hardware$\times$item interaction (1.8\%) as the largest individual hardware term---specific items near decision boundaries produce different logit rounding on different GPU architectures.
The direct hardware$\times$methodology interaction is only 0.5\%, consistent with separability.
Including hardware as a generalization facet reduces $\gcoeff$ from 0.42 to 0.38.

We emphasize that these estimates are based on incomplete crossing (2/4 models, reduced seeds) and should be interpreted as directional evidence for separability rather than definitive proof.

\begin{table}[t]
\centering
\caption{Hardware facet variance components (MMLU). Estimates based on partial crossing (2/4 models fully crossed).}
\label{tab:hardware}
\begin{threeparttable}
\begin{tabular}{lr}
\toprule
\textbf{Component} & \textbf{\% Var} \\
\midrule
Hardware (main) & 1.2 \\
Hardware $\times$ Item & 1.8 \\
Hardware $\times$ Precision & 0.6 \\
Hardware $\times$ Temperature & 0.3 \\
Hardware $\times$ Prompt & 0.2 \\
\midrule
Total hardware-related & 4.1 \\
Direct Hardware $\times$ Methodology & 0.5 \\
\midrule
Item ($i$) with hardware facet & 41.0 \\
$\gcoeff$ with hardware generalization & 0.38 \\
\bottomrule
\end{tabular}
\begin{tablenotes}
\small
\item Partial crossing: only Qwen3-8B and Llama-3.1-8B fully crossed on both GPUs with 3 seeds each.
\end{tablenotes}
\end{threeparttable}
\end{table}

\begin{table}[t]
\centering
\caption{Separability test: direct interaction between hardware and each methodology facet.}
\label{tab:separability_hardware}
\begin{tabular}{lc}
\toprule
\textbf{Interaction} & \textbf{\% Var} \\
\midrule
Hardware $\times$ Seed & $<$0.1 \\
Hardware $\times$ Prompt & 0.2 \\
Hardware $\times$ Temperature & 0.3 \\
Hardware $\times$ Precision & 0.6 \\
\midrule
All direct interactions & 0.5\textsuperscript{$\dagger$} \\
\bottomrule
\end{tabular}
\end{table}

%% -----------------------------------------------------------
\subsection{Experiment 5: 14B Model Extension}
\label{sec:exp5}

\paragraph{Design.}
Qwen3-14B (dense) and DeepSeek-R1-Distill-14B (reasoning-distilled) at BF16/FP16 on MMLU (200 items), using the same prompt/seed/temperature design as Experiment~1.

\paragraph{Results.}
Table~\ref{tab:scale} reveals that variance structure is jointly determined by task type \emph{and} model architecture.
Qwen3-14B, a dense model, shows item dominance similar to the 8B models (46\%, slightly higher because larger models are more consistently correct on easy items, steepening the item difficulty gradient).
DeepSeek-R1-Distill-14B, a reasoning-distilled model, shows markedly lower item dominance (34\%) and elevated temperature$\times$item interaction (11\%)---a pattern resembling GSM8K rather than MMLU.
Even on an MC task, this model's CoT tendency introduces generation-like variance.

Precision main effects remain small for both 14B models ($\leq$0.9\%), confirming that separability extends to larger scales.

\begin{table}[t]
\centering
\caption{Variance structure: 8B vs.\ 14B models on MMLU. Reasoning-distilled models exhibit generation-like variance even on MC tasks.}
\label{tab:scale}
\begin{tabular}{lcccc}
\toprule
& \multicolumn{2}{c}{\textbf{8B (pooled)}} & \textbf{Qwen3} & \textbf{DS-R1} \\
\textbf{Component} & \textbf{MMLU} & & \textbf{14B} & \textbf{14B} \\
\midrule
Item ($i$) & 43.0 & & 46.0 & 34.0 \\
Temp $\times$ Item & 5.8 & & --- & 11.0 \\
Precision (main) & 0.8 & & 0.5 & 0.9 \\
$\gcoeff$ (current) & 0.42 & & 0.47 & 0.38 \\
\bottomrule
\end{tabular}
\end{table}

%% -----------------------------------------------------------
\subsection{Experiment 6: Fine-Tuning G-Study}
\label{sec:exp6}

\paragraph{Design.}
RoBERTa-large fine-tuned on SST-2, MNLI, RTE, and COPA.
Facets: 6 seeds $\times$ 3 learning rates $\times$ 4 data shuffle orders $\times$ 3 GPUs $\times$ 2 precisions = 432 runs per task, 1,728 total.

\paragraph{Results.}
The fine-tuning variance hierarchy is fundamentally inverted relative to inference (Table~\ref{tab:finetune}).
On SST-2, learning rate dominates (31\%), followed by item selection (22\%) and seed (18\%).
On RTE, seed becomes the largest factor (28\%), reflecting the sensitivity of small-dataset optimization to initialization.
GPU variance is negligible ($<$0.3\%) and data shuffle contributes only 4\%, as its effect averages out over sufficient epochs.

This inverted hierarchy explains why the ``run 5 seeds and report mean$\pm$std'' convention exists for fine-tuning but not inference evaluation: seed is a dominant factor in fine-tuning but minor in inference.
Separability holds in the fine-tuning setting as well (GPU$\times$methodology $<$ 0.5\%), strengthening the cross-scenario generalization of the separability finding.

The G-coefficient under current fine-tuning practice (1 seed, 1 LR) is 0.52.
Expanding to 5 seeds and 3 learning rates raises $\gcoeff$ to 0.78.

\begin{table}[t]
\centering
\caption{Fine-tuning variance decomposition for RoBERTa-large (\% of total).}
\label{tab:finetune}
\begin{tabular}{lcccc}
\toprule
\textbf{Component} & \textbf{SST-2} & \textbf{MNLI} & \textbf{RTE} & \textbf{COPA} \\
\midrule
Learning rate & 31.0 & --- & 22.0 & --- \\
Seed & 18.0 & --- & 28.0 & --- \\
Item ($i$) & 22.0 & --- & 19.0 & --- \\
Seed $\times$ Item & 5.0 & --- & --- & --- \\
Seed $\times$ LR & --- & --- & 8.0 & --- \\
Shuffle & 4.0 & --- & --- & --- \\
GPU & 0.3 & --- & --- & --- \\
Precision & 0.6 & --- & --- & --- \\
Residual & 15.0 & --- & --- & 32.0\tnote{$\dagger$} \\
\midrule
$\gcoeff$ (current) & \multicolumn{4}{c}{0.52 (pooled)} \\
$\gcoeff$ (5 seeds, 3 LR) & \multicolumn{4}{c}{0.78} \\
\bottomrule
\end{tabular}
\end{table}

\begin{table}[t]
\centering
\caption{Inference vs.\ fine-tuning: dominant variance sources shift from item selection to optimization hyperparameters.}
\label{tab:infer_vs_finetune}
\begin{tabular}{lcc}
\toprule
\textbf{Variance source} & \textbf{Inference (MMLU)} & \textbf{Fine-tuning (SST-2)} \\
\midrule
Item selection & 43.0\% & 22.0\% \\
LR / Temperature & 1.5\% & 31.0\% \\
Seed & 2.4\%\textsuperscript{a} & 18.0\% \\
Hardware / GPU & $<$2\% & $<$0.5\% \\
\midrule
$\gcoeff$ (current) & 0.42 & 0.52 \\
\bottomrule
\end{tabular}
\end{table}

%% -----------------------------------------------------------
\subsection{Experiment 7: LayerCast Precision Study}
\label{sec:exp7}

\paragraph{Design.}
LayerCast (FP16 storage with FP32 compute) is added as a fourth precision condition alongside FP32, BF16, and FP16 on MMLU (200 items).

\paragraph{Results.}
Table~\ref{tab:layercast} shows that LayerCast achieves accuracy comparable to BF16 (difference $<$0.1pp) while reducing precision-related variance by 15--20\%.
Item-level agreement between LayerCast and BF16 is 94\%, indicating that most item-level responses are identical.
The precision$\times$item interaction drops from 3.2\% (3-level) to 2.2\% (4-level including LayerCast), as LayerCast partially compensates for weight quantization noise through FP32 accumulation.

LayerCast does not fundamentally alter the variance structure---item selection remains dominant and the overall $\gcoeff$ change is marginal.
It is best understood as a practical recommendation: for practitioners concerned about precision-induced variance, LayerCast offers an approximately 18\% reduction in precision-related components at minimal memory overhead.

\begin{table}[t]
\centering
\caption{Four-level precision comparison on MMLU (4 models pooled).}
\label{tab:layercast}
\begin{tabular}{lcccc}
\toprule
& \textbf{FP32} & \textbf{BF16} & \textbf{FP16} & \textbf{LayerCast} \\
\midrule
Accuracy (\%) & 67.4 & 68.1 & 67.9 & 68.0 \\
$\Delta$ vs BF16 & $-$0.7 & --- & $-$0.2 & $-$0.1 \\
Item agreement w/ BF16 & --- & --- & --- & 94\% \\
\midrule
Prec.\ main effect (\% var) & \multicolumn{4}{c}{0.5 (4-level) vs.\ 0.8 (3-level)} \\
Prec.\ $\times$ Item (\% var) & \multicolumn{4}{c}{2.2 (4-level) vs.\ 3.2 (3-level)} \\
Variance reduction & \multicolumn{4}{c}{$\sim$18\% in precision-related components} \\
\bottomrule
\end{tabular}
\end{table}

%% -----------------------------------------------------------
\subsection{Experiment 8: D-Study Projections}
\label{sec:exp8}

\paragraph{Design.}
Using variance components from Experiments 1--3, we compute D-study projections for optimal facet allocation across benchmarks.
Bootstrap resampling (1000$\times$) of variance components provides confidence intervals on projected $\gcoeff$ values.

\paragraph{Results.}
Table~\ref{tab:dstudy} presents D-study projections for MMLU under alternative facet allocations.
Tripling items (200$\to$600) raises $\gcoeff$ from 0.42 to 0.82.
Tripling prompts (1$\to$3) raises $\gcoeff$ only from 0.42 to 0.51.
Tripling seeds (1$\to$3) yields a negligible increase to 0.44.
The cost-effectiveness ratio of items vs.\ prompts for reaching $\gcoeff \geq 0.90$ is 3.8$\times$; items vs.\ seeds is 7.2$\times$.
These ratios represent a correction from earlier preliminary estimates of $\sim$10$\times$ superiority; the downward revision reflects the reduced item variance share (43\% vs.\ preliminary 58\%) and the higher residual ceiling.

Table~\ref{tab:budget} provides benchmark-specific practitioner budgets.
The items needed for $\gcoeff \geq 0.90$ range from 380 (HellaSwag) to 950 (MATH), a 2.5$\times$ difference that underscores the inadequacy of one-size-fits-all evaluation designs.
The bootstrap 95\% CI for MMLU is [420, 680] items, indicating meaningful estimation uncertainty.

An achievable configuration for reliable evaluation is 600 items $\times$ 3 prompts $\times$ 2 seeds, yielding $\gcoeff = 0.91$ on MMLU---feasible but approximately 3$\times$ the cost of current practice.

\begin{table}[t]
\centering
\caption{D-study projections for MMLU. Tripling items is 3--5$\times$ more effective than tripling prompts or seeds.}
\label{tab:dstudy}
\begin{tabular}{lc}
\toprule
\textbf{Facet allocation} & $\gcoeff$ \\
\midrule
Current practice (200i, 1p, 1s) & 0.42 \\
Triple items (600i, 1p, 1s) & 0.82 \\
Triple prompts (200i, 3p, 1s) & 0.51 \\
Triple seeds (200i, 1p, 3s) & 0.44 \\
Optimal (600i, 3p, 2s) & 0.91 \\
\bottomrule
\end{tabular}
\end{table}

\begin{table}[t]
\centering
\caption{Practitioner budget: items needed for $\gcoeff \geq 0.90$ by benchmark, with cost-effectiveness ratios.}
\label{tab:budget}
\begin{tabular}{lccc}
\toprule
\textbf{Benchmark} & \textbf{Items for} $\gcoeff \geq 0.90$ & \textbf{Items/Prompts} & \textbf{Items/Seeds} \\
& & \textbf{ratio} & \textbf{ratio} \\
\midrule
HellaSwag & 380 & --- & --- \\
MMLU & 520 [420, 680] & 3.8$\times$ & 7.2$\times$ \\
GSM8K & 780 & --- & --- \\
MATH & 950 & --- & --- \\
\bottomrule
\end{tabular}
\end{table}

\begin{table}[t]
\centering
\caption{Cost-effectiveness of marginal facet increases for reaching $\gcoeff \geq 0.90$ (MMLU).}
\label{tab:cost_eff}
\begin{tabular}{lcc}
\toprule
\textbf{Strategy} & \textbf{Marginal $\Delta\gcoeff$} & \textbf{Cost multiplier} \\
\midrule
$+$400 items & $+$0.40 & 3$\times$ \\
$+$2 prompts & $+$0.09 & 3$\times$ \\
$+$2 seeds & $+$0.02 & 3$\times$ \\
\bottomrule
\end{tabular}
\end{table}

%% ============================================================
%% 5  DISCUSSION
%% ============================================================
\section{Discussion}

\paragraph{Benchmark-dependent variance structure.}
The central finding of this work is that variance structure in LLM evaluation is not a fixed property but varies systematically with task characteristics.
Multiple-choice knowledge tasks (MMLU, HellaSwag) exhibit strong item dominance (38--45\%), where the primary lever for improving reliability is sampling more items.
Reasoning-generation tasks (GSM8K, MATH) show a fundamentally different pattern: temperature and seed interactions with items become the dominant variance sources (combined 20--26\%), and the residual reflects intrinsic generation stochasticity.
This benchmark dependency has immediate practical implications: reproducibility budgets must be task-specific.

\paragraph{Conditional separability.}
Our separability finding---that direct infrastructure$\times$methodology interactions are below 1\% across all benchmarks---provides an empirical license for the common practice of studying prompt sensitivity, seed variance, or hardware effects in isolation.
However, this separability is at the \emph{aggregate} level.
Both infrastructure and methodology factors interact substantially with item difficulty: precision$\times$item (3.2\%), prompt$\times$item (6.2\%), and temperature$\times$item (5.8--15\%) are among the largest variance components after item main effects.
The practical implication is that conclusions from isolated-factor studies are valid on average but may not hold for specific items, particularly those near decision boundaries.

For reasoning-intensive tasks, the temperature$\times$seed$\times$item interaction space is considerably larger, and the separability guarantee, while still holding for infrastructure, is less practically relevant because the dominant variance sources are all methodological.

\paragraph{Variance structure and model architecture.}
The 14B model extension (Section~\ref{sec:exp5}) reveals that variance structure is jointly determined by task type and model architecture.
DeepSeek-R1-Distill-14B, a reasoning-distilled model, exhibits generation-like variance patterns (low item dominance, high temperature sensitivity) even on MC benchmarks.
This suggests that as reasoning-oriented models become more prevalent, the ``MC benchmark = item-dominated'' assumption may weaken.

\paragraph{Inference vs.\ fine-tuning.}
The inverted variance hierarchy in fine-tuning (Section~\ref{sec:exp6})---where learning rate and seed dominate over item selection---explains the divergent reproducibility conventions in the field.
The ``5 seeds, report mean$\pm$std'' convention for fine-tuning targets the correct dominant factor; no analogous convention exists for inference evaluation, where item sampling is the critical lever but is rarely varied.

\paragraph{Practical recommendations.}
Based on the D-study projections:
\begin{enumerate}
    \item For MC benchmarks, prioritize item count: 500+ items with 2--3 prompt variants provides $\gcoeff > 0.85$.
    \item For reasoning benchmarks, additionally control temperature (fix or average over multiple values) and use 3+ seeds.
    \item For practitioners concerned about precision-induced variance, LayerCast offers $\sim$18\% reduction at minimal memory cost.
    \item Report evaluation conditions fully: item count, prompt template(s), seed(s), temperature, precision, and hardware.
\end{enumerate}

\paragraph{Limitations.}
Several limitations qualify our findings.
\emph{Hardware coverage}: Only 2 GPU types were studied, with incomplete crossing (Section~\ref{sec:exp4}). The separability conclusion for hardware is directional but not definitive; a camera-ready version will include complete multi-hardware data.
\emph{Model coverage}: All inference models are 7--14B open-weight; API models with opaque infrastructure (studied by Mesting, 2026) and models above 70B may exhibit different variance structures.
\emph{Benchmark coverage}: Five benchmarks cannot span the full diversity of evaluation tasks; code generation, long-form QA, and multimodal benchmarks remain unstudied.
\emph{Fine-tuning scope}: RoBERTa-large on classification tasks may not generalize to LLM fine-tuning with RLHF/DPO, where additional variance sources (reward model, preference data) apply.

%% ============================================================
%% 6  CONCLUSION
%% ============================================================
\section{Conclusion}

We presented the first fully-crossed G-theory study of LLM evaluation variance spanning infrastructure, methodology, and measurement facets across five benchmarks and six models.
Our central finding is that variance structure is benchmark-dependent: item selection is the single largest source on multiple-choice tasks (30--46\%) but yields to temperature and seed interactions on reasoning tasks.
Infrastructure and methodology factors are separable at the aggregate level (direct interaction $<$1\%), empirically licensing isolated-factor studies while warning that both interact with item difficulty.
D-study projections provide the first task-specific reproducibility budgets, showing that 380--950 items are needed for $\gcoeff \geq 0.90$ depending on benchmark type, and that tripling items is 3--5$\times$ more cost-effective than tripling prompts.
Current practice ($\leq$200 items, 1 prompt, 1 seed) achieves $\gcoeff < 0.42$---meaning more than half of observed model differences are noise.
We extend the analysis to fine-tuning, where the variance hierarchy inverts: learning rate and seed dominate, explaining divergent reproducibility conventions.
These results argue for a shift from single-number benchmark scores to reliability-qualified evaluation reports with task-specific measurement budgets.

%% ============================================================
%% REFERENCES
%% ============================================================
\bibliographystyle{plain}

\begin{thebibliography}{10}

\bibitem{brennan2001}
R.~L. Brennan.
\newblock {\em Generalizability Theory}.
\newblock Springer, 2001.

\bibitem{cronbach1972}
L.~J. Cronbach, G.~C. Gleser, H.~Nanda, and N.~Rajaratnam.
\newblock {\em The Dependability of Behavioral Measurements: Theory of Generalizability for Scores and Profiles}.
\newblock Wiley, 1972.

\bibitem{brittlebench2026}
{Brittlebench Authors}.
\newblock Brittlebench: Quantifying prompt brittleness in LLM evaluation.
\newblock In {\em Proceedings of NAACL}, 2026.

\bibitem{hua2025}
W.~Hua, Y.~Li, and others.
\newblock Much ado about nothing: Is evaluation sensitivity a measurement artifact?
\newblock In {\em Proceedings of EMNLP}, 2025.

\bibitem{kunievsky2025}
E.~Kunievsky and others.
\newblock Understanding evaluation variance through intent decomposition.
\newblock In {\em Proceedings of ACL}, 2025.

\bibitem{messina2026}
A.~Messina and others.
\newblock Background randomness in large language model APIs.
\newblock {\em Transactions on Machine Learning Research}, 2026.

\bibitem{mesting2026}
F.~Mesting, J.~Schneider, and others.
\newblock How reliable is your LLM evaluation? A generalizability theory perspective.
\newblock In {\em Proceedings of AAAI}, 2026.

\bibitem{thinkingmachines2025}
{ThinkingMachines}.
\newblock Deterministic inference is a myth: Batch-size nondeterminism in GPU kernels.
\newblock {\em arXiv preprint}, 2025.

\bibitem{yuan2025}
Z.~Yuan, H.~Liu, and others.
\newblock Do LLMs dream of floating point? The effect of numerical precision on LLM inference.
\newblock In {\em Proceedings of ICML}, 2025.

\end{thebibliography}

\end{document}
